\hypertarget{appendix-f2ee4d086415}{}\label{referencestart:appendix-f2ee4d086415}
\section{Prefix scan and stream compaction}\label{appendix-f2ee4d086415-prefix-scan-and-stream-compaction}

\hypertarget{appendix-f2ee4d086415-prefix-scan-and-stream-compaction}{}

A prefix scan turns local counts into globally consistent offsets. For an exclusive sum scan of input \(x\),

\[
y_0=0,
\qquad
y_i=\sum_{j=0}^{i-1}x_j.
\]

If each element says how many output records it will emit, \(y_i\) is the first output slot owned by element \(i\). This converts a race-prone global append into deterministic non-overlapping regions.

For Boolean compaction, let predicate \(b_i\in\{0,1\}\). Scan \(b\) to obtain offsets. Every active element with \(b_i=1\) writes itself to output{[}scan{[}i{]}{]}. The final output count is scan{[}n-1{]} + b{[}n-1{]}.

\subsection{Parallel structure}\label{appendix-f2ee4d086415-parallel-structure}

\hypertarget{appendix-f2ee4d086415-parallel-structure}{}

The naive parallel scan that repeatedly doubles an offset performs \(O(n\log n)\) work. A work-efficient tree scan performs an up-sweep/reduction followed by a down-sweep and uses \(O(n)\) work with \(O(\log n)\) depth.

On a GPU, a practical large-array scan normally has several levels:

\begin{enumerate}
\def\labelenumi{\arabic{enumi}.}
\tightlist
\item
  scan each workgroup block into local output;
\item
  record each block total;
\item
  scan the block totals;
\item
  add each scanned block offset back into that block's local results.
\end{enumerate}

Every phase has a dependency. Dispatching all phases without establishing visibility/order is not a valid scan merely because each shader is individually correct.

Shared-memory bank conflicts and subgroup width affect performance, while workgroup barriers only synchronize participating invocations in that workgroup.

\subsection{Stream compaction as an allocation proof}\label{appendix-f2ee4d086415-stream-compaction-as-an-allocation-proof}

\hypertarget{appendix-f2ee4d086415-stream-compaction-as-an-allocation-proof}{}

Compaction is useful beyond filtering visible objects. It appears in GPU-driven draw generation, BVH construction, particle survival, tiled light lists, Gaussian tile duplication and sparse work queues.

The scan proves that each surviving producer owns a distinct range before any scatter writes occur.

\subsection{Evidence boundary}\label{appendix-f2ee4d086415-evidence-boundary}

\hypertarget{appendix-f2ee4d086415-evidence-boundary}{}

The current repository has CPU-side scan/compaction reasoning in the canonical volumes but does not yet claim a native Vulkan/Metal scan kernel benchmark. The next hardware proof should compare a scalar CPU oracle, a workgroup implementation, multi-workgroup hierarchy, adversarial non-power-of-two sizes, and captured GPU timing.

Primary background: NVIDIA GPU Gems 3, Chapter 39, ``Parallel Prefix Sum (Scan) with CUDA.''

Teach and reconstruct: \href{../06-GPU-Architecture-and-Real-Time-Rendering-P01/main.pdf\#01-foundations-i-ch008-scan-turns-local-counts-into-nonoverlapping-output-regions}{Continuous Notes/Computing Foundations - From Bits to C++ and Systems APIs\#Scan turns local counts into nonoverlapping output regions (GPU Architecture and Real-Time Rendering, Part 1, Ch. 1)}.

Related: \href{../06-GPU-Architecture-and-Real-Time-Rendering-P01/main.pdf\#appendix-f662b5b1773b}{GPU workgroups and synchronization (GPU Architecture and Real-Time Rendering, Part 1, Reference notes: GPU workgroups and synchronization)}, \href{../06-GPU-Architecture-and-Real-Time-Rendering-P02/main.pdf\#appendix-5653a7aea6bb}{BVH build traversal and workload selection (GPU Architecture and Real-Time Rendering, Part 2, Reference notes: BVH build traversal and workload selection)}.
